% Options for packages loaded elsewhere
\PassOptionsToPackage{unicode}{hyperref}
\PassOptionsToPackage{hyphens}{url}
%
\documentclass[
]{article}
\usepackage{amsmath,amssymb}
\usepackage{iftex}
\ifPDFTeX
  \usepackage[T1]{fontenc}
  \usepackage[utf8]{inputenc}
  \usepackage{textcomp} % provide euro and other symbols
\else % if luatex or xetex
  \usepackage{unicode-math} % this also loads fontspec
  \defaultfontfeatures{Scale=MatchLowercase}
  \defaultfontfeatures[\rmfamily]{Ligatures=TeX,Scale=1}
\fi
\usepackage{lmodern}
\ifPDFTeX\else
  % xetex/luatex font selection
\fi
% Use upquote if available, for straight quotes in verbatim environments
\IfFileExists{upquote.sty}{\usepackage{upquote}}{}
\IfFileExists{microtype.sty}{% use microtype if available
  \usepackage[]{microtype}
  \UseMicrotypeSet[protrusion]{basicmath} % disable protrusion for tt fonts
}{}
\makeatletter
\@ifundefined{KOMAClassName}{% if non-KOMA class
  \IfFileExists{parskip.sty}{%
    \usepackage{parskip}
  }{% else
    \setlength{\parindent}{0pt}
    \setlength{\parskip}{6pt plus 2pt minus 1pt}}
}{% if KOMA class
  \KOMAoptions{parskip=half}}
\makeatother
\usepackage{xcolor}
\usepackage{longtable,booktabs,array}
\usepackage{calc} % for calculating minipage widths
% Correct order of tables after \paragraph or \subparagraph
\usepackage{etoolbox}
\makeatletter
\patchcmd\longtable{\par}{\if@noskipsec\mbox{}\fi\par}{}{}
\makeatother
% Allow footnotes in longtable head/foot
\IfFileExists{footnotehyper.sty}{\usepackage{footnotehyper}}{\usepackage{footnote}}
\makesavenoteenv{longtable}
\usepackage{graphicx}
\makeatletter
\def\maxwidth{\ifdim\Gin@nat@width>\linewidth\linewidth\else\Gin@nat@width\fi}
\def\maxheight{\ifdim\Gin@nat@height>\textheight\textheight\else\Gin@nat@height\fi}
\makeatother
% Scale images if necessary, so that they will not overflow the page
% margins by default, and it is still possible to overwrite the defaults
% using explicit options in \includegraphics[width, height, ...]{}
\setkeys{Gin}{width=\maxwidth,height=\maxheight,keepaspectratio}
% Set default figure placement to htbp
\makeatletter
\def\fps@figure{htbp}
\makeatother
\setlength{\emergencystretch}{3em} % prevent overfull lines
\providecommand{\tightlist}{%
  \setlength{\itemsep}{0pt}\setlength{\parskip}{0pt}}
\setcounter{secnumdepth}{-\maxdimen} % remove section numbering
\ifLuaTeX
  \usepackage{selnolig}  % disable illegal ligatures
\fi
\usepackage{bookmark}
\IfFileExists{xurl.sty}{\usepackage{xurl}}{} % add URL line breaks if available
\urlstyle{same}
\hypersetup{
  pdftitle={Culture consumption and social capital},
  hidelinks,
  pdfcreator={LaTeX via pandoc}}

\title{Culture consumption and social capital\footnote{\textsuperscript{*}Data
  for this project were obtained from the Association of Religion Data
  Archives at \url{http://www.thearda.com}. The author bears sole
  responsibility for the current interpretation, analyses and
  tabulations of these data. I would like to thank Paul DiMaggio, Bart
  Bonikowski, and the members of the Princeton seminar in cultural
  sociology for their very useful comments on previous drafts of this
  paper.}}
\author{}
\date{}

\begin{document}
\maketitle

Omar Lizardo

Department of Sociology

University of Notre Dame

810 Flanner Hall

Notre Dame, IN, 46556

\href{mailto:olizardo@nd.edu}{\nolinkurl{olizardo@nd.edu}}

Words: 12,254

Last Revised: 08-Feb-24

\subsubsection{Abstract}\label{abstract}

\subsection{\texorpdfstring{\hfill\break
Introduction}{ Introduction}}\label{introduction}

Prevailing models of taste formation, transmission and decay in the
sociology of culture posit that current tastes are largely a product of
the individual's current network configuration. The standard assumption
is that tastes are transmitted through existing network ties, in
particular those social connections that are characterized by a high
degree of sociodemographic similarity---``status homophily'' (Lazarsfeld
and Merton 1954)---between ego and alter (McPherson, Smith-Lovin, and
Cook 2001). From this perspective, birds of a feather not only flock
together, but largely due to this mutual flocking they end up singing,
going to the movies, and reading together (Mark 1998b, 2003). A key
explanatory strength of the network model is its ability to account for
the fact that the distribution of tastes across persons in
socio-demographic space is ``lumpy'' and clustered in proximal regions
of this space. The model explains this by suggesting that the dynamics
of cultural taste formation and transmission is best envisioned as a
``local bandwagon'' process of network-based imitation (Mark 2003).
Individuals are construed as highly susceptible to the influence of
their current contacts and by implication tastes are seen not as
enduring aspects of the person but as malleable, ``thin'' contents
liable to change when those network ties in turn change in their turn.

And change they do. Recent longitudinal network research has shown that
rather than being stable over time personal networks as constantly
shifting with new contacts constantly being added and deleted throughout
the adult life-course (Burt 2000, 2002; Morgan, Neal, and Carde 1997;
Suitor, Wellman, and Morgan 1997; Wellman, Wong, Tindall, and Nazer
1997). Thus, what is truly surprising about over-time network
``stability'' over time is ``how \emph{unstable} intimacy is'' (Wellman
\emph{et al.} 1997: 47). The empirical fact of network volatility
coupled with the network assumption of taste as being largely a function
of the cultural choices of the individual's current set of contacts,
implies large-scale volatility of those very same cultural tastes. Thus
if we take the network perspective seriously, we should expect that as
friends change, so do tastes, so we should expect small to moderate
correlations between an individual's taste at time \emph{t} and her
taste at time \emph{t+n}.

Very different assumptions are at core of the cultural capital model of
taste formation and acquisition (Bourdieu 1984; Holt 1998; Lizardo
2006). Instead of positing tastes as highly malleable contents at the
perennial mercy of a constantly shifting network configuration, cultural
tastes are conceived as highly \emph{durable} dispositions toward broad
types of cultural objects and performances (Frith 1998). These
dispositions are formed in the family and school environment throughout
childhood and early to middle adolescence (Smith 1995). Once set, they
are highly resistant to change. The reason for this is that as Holt
(1997) notes following Bourdieu (1990), tastes are part the person's
\emph{embodied cultural capital} and thus represents a highly resilient
aspect of the individual's relationship to cultural objects (Bourdieu
1984), encoded in the memory traces, semi-conscious affective tendencies
and habitual dispositions responsible for practical action (Bourdieu
2000).

Contrary to the network account, according to the cultural capital model
tastes should not exhibit over-time volatility, but instead show high
inter-temporal correlation. Lizardo (2006) uses this stability
assumption to suggest that tastes can in fact be seen as probable
(reciprocal) \emph{drivers} of network connections rather than being
seen as exclusively \emph{determined} by those connections. Using
two-stage instrumental variables and matching methods on cross-sectional
data, that study uncovers evidence consistent with this reciprocal
effect. Engagement in a wide variety of cultural practices increases
network size net of standard controls and after statistically correcting
for simultaneity bias. This is in agreement with dynamic theories of
cultural and network evolution, which rather than postulating a one-way
avenue of causation going from networks to cultural tastes propose a
reciprocal (over-time) two-way casual interaction between the two
(Carley 1991; Mark 1998a). However, in the aforementioned study the
stability hypothesis remains an empirically unverified background
assumption. Furthermore to reliance on cross-sectional data, the casual
effect of cultural taste on network maintenance or even the more widely
accepted effect of network turnover on cultural taste loss, has not been
decisively established.

In this paper, I use individual-level longitudinal panel data to (1)
empirically adjudicate between these divergent empirical scenarios
regarding stability implied by the network and cultural capital models,
and to establish (2) direct and unambiguous casual connections between
network turnover and cultural taste loss and (3) high rates of cultural
activity and the maintenance of high rates of social connectivity in a
longitudinal data analysis framework. In the next section, I go on to
develop the empirical implications of the network and cultural capital
views on more detail. In the following section, I introduce the data and
propose my analytic strategy. In the next section, I present the
results. I close with a discussion section highlighting the implications
of the results for current theoretical understanding of the relationship
between cultural tastes and social networks.

\subsection{Network versus cultural capital theories of cultural taste
loss: some empirical
implications}\label{network-versus-cultural-capital-theories-of-cultural-taste-loss-some-empirical-implications}

\subsubsection{Network theory}\label{network-theory}

Contemporary network theories of cultural taste and attitude formation
and change (and by implication decay) suggest that the contemporary
taste profile of an individual is largely a function of the tastes held
by the other persons in her surrounding social network (Erickson 1982,
1988). The key idea is that social ties serves as the ``conduits''
through which new cultural practices are \emph{transmitted} and by way
of which individuals \emph{learn} these new cultural practices from
their social contacts (Carley 1991; Mark 1998a). This proposition is
shared by both straightforward network accounts (i.e. Cardon and Granjon
2005; Erickson 1996; Kane 2004) and more ambitious ecological-network
theories of cultural taste (Mark 1998a, 2003; McPherson 2004). Mark
(2003: 323-324) for instance speaks of social ties to friends and family
as the most important sources of cultural taste acquisition:
``{[}p{]}eople's friends and family\ldots influence the leisure
activities they participate in\ldots participation in various sports,
visiting art museums, and gardening are\ldots leisure activities that
people often lean from parents, siblings, and/or friends.''

Network theories of taste transmission implicitly presume a relatively
stable ``social structure'' coupled to relatively unstable patterns of
cultural taste, otherwise it would be hard to think how something that
varies at a more rapid pace ``causes'' or ``affects'' something that is
more stable than itself (Simon 1952). It bears mention that this
relative stability assumption is shared by the entire family of network
theories that conceive of social structure as the pipes through which
contents, or information ``flows'' and diffuses from one actor to
another (Borgatti 2005). This implicit (and crucial) assumption of the
stability of social structure, when this term is defined---as it is
usually done in network theory---as the concrete patterns of network
relations that bind individuals to one another (Berkowitz 1982, 1988;
Wellman 1983, 1988; Wellman and Berkowitz 1988), can certainly be
questioned in light of recent longitudinal research on constancy of
network relations over time. This research---see the 1997 and 2005
special issues of \emph{Social Networks} (volume 19, number 1 and volume
27 number 4) and Burt 2000, for a thorough review)---shows that instead
of being constant, network relations are extremely volatile, even when
the time frame is constrained to relatively short periods of time, such
as a year or less (Bidart and Degenne 2005; Suitor, Wellman, and Morgan
1997).

Most longitudinal studies of personal networks reveal surprisingly high
levels of instability in personal networks over time (Lang 2000), with
anywhere approximately about 30 to 60\% of the personal network
experiencing turnover in about a year's time. In his own research on
bankers, Burt (2000: 11) finds that ``three in four of the 22,709
colleague relations at risk of decay\ldots are gone next year.'' This
means that old ties are constantly being deleted---what Burt (2000)
refers to as ``tie decay''---and new ones formed constantly throughout
the life-course (Suitor, Wellman, and Morgan 1997; Wellman \emph{et al.}
1997). Suitor and Keeton (1997: 57), report that more than half of the
associates...named at T1 were not named again at either T2 or T3.''
Morgan et al (1997: 14) conclude that their results ``point to a
considerable degree of instability in who is present in the network over
time. In particular, the average overlap rate of 55\% means that just
over half of those who were named in one interview or the other appeared
in both. Similarly, interpreting the average phi value of 0.34 as a
test-retest correlation indicates only a limited ability to predict who
was in the network from one time to the next.''

These findings imply that for most individuals a recurrent process of
new contact selection, extant relationship ``fine tuning'' and possible
deletion of old social ties is constantly at work (Lazarsfeld and Merton
1954). This represents an important---if ill understood---facet
underlying the temporal evolution of a person's relational
micro-environment. This dynamic acquires renewed salience in light of
the observation that recently formed network connections appear to
exhibit what Burt (2000), drawing on Stinchcombe (1965: 148-150), refers
to as a ``liability of newness.'' This means that dyadic network ties
tend to display a high probability of dissolution at early stages of the
relationship, with the probability of termination steeply declining as
the relationship matures. This proposition implies that the ``tie decay
function'' (that is the probability of a tie being dissolved) is
non-linearly decreasing in relationship age, with a high probability for
young relationships which decreases as these dyadic ties ripen into
close friendships. Burt (2000, 2002) finds evidence consistent with this
hypothesis for both strong ties and network bridges (weak ties),
respectively.

It is true that network stability is much more substantial in the
``core'' personal network, with the bulk of the instability concentrated
on the periphery (Bidart and Lavenu 2005; Suitor, Wellman, and Morgan
1997; Wellman \emph{et al.} 1997). However, most longitudinal network
studies show that the stable ``core'' network tends to be composed of a
small set of alters, primarily kin. The periphery of the network
includes most of the persons usually referred to as (non-kin)
``friends.'' Studies show that there is considerable turnover even among
those ties usually considered ``close friends'' (Bidart and Lavenu
2005). It is possible to revise network theories of taste transmission
to suggest that taste are mostly transmitted through the small,
relatively stable core and not through the larger, unstable periphery.
However, most contemporary network theories of taste do not make this
distinction, and assume that taste are as equally likely to be learned
from family and friends (Mark 2003). It is clear from recent studies on
personal communities that the ``core'' kin network while mobilized for
many purposes, including emotional and financial support, is not used
for companionship and sociability (Simmel 1949) in any meaningful sense.
Most research and theory on culture consumption communities shows that
if tastes are exploited in interaction, this is done precisely with
those contacts that are usually named as ``friends'' and which are seen
outside the home and the neighborhood (DiMaggio 1987; Fiske 1987; Frith
1998). However, there is nothing inherent in the logic of the network
theory of taste transmission that makes it incompatible with current
data suggesting large rates of network turnover and the liability of
newness of social ties.

Nevertheless, it is important to note that the social transmission
hypothesis coupled with the empirical observation of the instability of
social ties, jointly imply the inter-temporal instability of cultural
tastes. The reason for this is that as Mark (2003) notes a taste that is
not ``reinforced'' by social network influence through interaction is
likely to be ``forgotten.'' This is a key assumption of the network
based model of taste acquisition. Thus, if we retain the network
assumption that individual tastes are largely a function of the tastes
displayed by the persons current network alters, and if we accept the
fact that these alters are constantly change over the adult life course,
then we should expect cultural tastes to themselves exhibit high
volatility over time:

\emph{H1: We should observe low to moderate inter-temporal stability of
cultural taste across individuals.}

Network theory of course, does not only imply volatile tastes, but
directly suggests that any observed amount of taste loss should be
primarily associated with life-course events associated with network
turnover (i.e. geographical mobility, marriage and divorce, entering and
leaving the workforce, adding or dropping social interaction foci,
etc.). Surprisingly, but primarily due to the paucity of longitudinal
data of cultural taste, this most crucial of assumptions of the
network-theoretic view of cultural taste loss has never been assessed in
an unambiguous way, although cross-sectional analyses appear to be
consistent with computer simulated predictive scenarios (Carley 1991;
Mark 2003). The network hypothesis therefore implies that:

\emph{H2: Observed cultural taste loss should be primarily a function of
life-course changes associated with network turnover.}

\subsubsection{Cultural capital theory}\label{cultural-capital-theory}

The cultural capital approach to taste formation provides a very
different picture of the stability of tastes than that which can be
gleaned from contemporary network theory. Instead of implying that
tastes are unstable, ``shallow'' contents liable to be replaced at the
same pace that network relations are replaced, the cultural capital
approach suggest that taste are ``deep'' \emph{durable dispositions and
trans-temporally stable forms of embodied cultural capital} (Bourdieu
1986; Holt 1998). It is the relative stability of tastes vis a vis
network relations that allows them to function as resources to help
\emph{sustain} and \emph{maintain} these network relations over
time---primarily by serving as topical resources in conversational
interaction rituals (DiMaggio 1987; Lizardo 2006) providing a common
focus of attention and affective commitment (Collins 2004). This helps
to shift interactions away from purely ``instrumental'' purposes of
information gathering and rational goal attainment to sociable
interaction for its own sake (Granovetter 2002; Simmel 1949).

While there has been little empirical research on the dynamic stability
of tastes through time, there is good reason to suppose that they are
more stable than most network theory leads us to believe. For instance,
most studies of the role of early family and school experiences on
cultural capital have shown strong effects of arts participation,
education and after-school training during adolescence on adult tastes
even after controlling for subsequent educational attainment (Kracman
1996). Smith (1995) using data from the 1993 General Social Survey,
shows that musical tastes are developed early in young adulthood and
appear to be fairly stable across age-cohorts. Holbrook and Schindler
(1989) find that persons tend to like styles of popular music with which
they were familiar in late-adolescence over older or more contemporary
styles, suggesting that tastes tend to be retained over the life-course.
Bennett, Emisson and Frow (1999: 178), using qualitative evidence from
retrospective interview reports on musical preferences in Australia find
that ``\ldots many of those who were in middle age had retained the
music preference of their late teens and early adulthood.'' They
conclude that while they ``\ldots cannot discount the possibility that
some people will undergo radical changes in their tastes as they age, on
balance the evidence appears to favor the explanation that the taste
acquired in the formative years of a cohort persist into later life''
(1999: 180). Crook (1997) finds strong correlations of reading and
\emph{beaux arts} practices on a three-year panel survey of the
Australian population and also across self-reports of cultural activity
in adolescence and adulthood, suggesting high stability of cultural
practices across the life-course. From this perspective therefore, while
the effect of social network change on cultural taste loss is not
denied, whatever relational (or ``network'') effects on taste exist,
must operate in the background of high-levels of cross-temporal taste
stability and the reciprocal action of cultural taste on social
networks:

\emph{H3: We should observe high-levels inter-temporal stability of
cultural taste across individuals.}

According to the cultural-capital theory inspired (i.e. Bourdieu 1986)
``conversion model'' of cultural into social capital (DiMaggio and Mohr
1985; Dimaggio and Useem 1978155; Lizardo 2006), the primary ways in
which tastes are used to sustain social networks, is by allowing the
individual to connect to alters that are both relationally close---when
specialist or ``niche'' cultural tastes are concerned---and relationally
distant---when talking about aptitude for ``popular'' cultural forms.
From this perspective, culture consumption and the ability to decode and
appropriate symbolic goods produced in the urban, folk and
``culture-industry'' realms (Crane 1993) are important resources that
individuals mobilize in the \emph{construction} and \emph{maintenance}
of active social relationships (Cardon and Granjon 2005; DiMaggio 1987).
In this respect the culture conversion model is consistent with
Lazarsfeld and Merton's (1954) process model of friendship formation,
with similarity based on cultural taste and practices (DiMaggio 1987;
Mark 2003) playing the role that ``value-homophily'' played in the
original formulation.\footnote{Contemporary network models of taste
  formation through social network ties (Mark 1998b) can be seen as
  partially based on the Lazarsfeld-Merton formulation. However, these
  models raise what was explicitly a secondary dynamic in the original
  scheme---i.e. mutual ``value-adjustment'' through social interaction
  of initially different value dispositions (Lazarsfeld and Merton 1954:
  33)---to being the only operative, and thus primary dynamic. It is
  clear however that for Lazarsfeld and Merton ``selective'' processes
  whereby contacts are sifted based on similarity of attitudinal
  dispositions similar to those emphasized in the culture-conversion
  account were considered primary (1954: 22-29).} As Cardon and Granjon
(2005: 302) note in their recent study of the role of culture
consumption in the formation of social connections among young adults in
France, ``\ldots a large proportion of ordinary conversational
interaction is generally devoted to comments\ldots on cultural practices
and recreational consumption. These are recurrent topics, present in a
wide range of social situations\ldots Many social relations are
constructed, in different ways, around cultural and recreational
activities.''

That the origins of new---and reproduction of previously
existing---network connections comes from the deployment of cultural
knowledge in interaction is implicit in previous models of the
relationship between culture and social structure (Fine 1979). For
instance, the Carley-Mark (Carley 1991, 1995; Mark 1998a, 2003)
``constructural'' model can be interpreted as an elementary schema of
how culture can be translated into social connections, and how social
structure (the distribution of chances to interact across persons in the
system) and cultural structure (the distribution of cultural forms
across persons) can be defined in an interdependent manner. The
constructural model is based on a simple assumption of similarity
(homophily), which postulates that the likelihood of a social tie
increases with the cultural similarity between any given dyad (a dynamic
process similar to the one proposed in Homans 1950 and Lazarsfeld and
Merton 1954). In this way the probability that two persons will interact
is driven by their cultural similarity. Interaction, in this positive
feedback loop, in its turn increases cultural similarity, as individuals
exchange their stocks of knowledge with one another.

Thus we should expect that individuals who command a wide variety of
tastes which cut across these well defined boundaries (i.e. those who
have been referred to as ``omnivores'' in recent research {[}Peterson
1992, 1997, 2005{]}) and who, given the stability assumption, will also
be able to ``keep'' that ``tolerant'' taste pattern active over
time---will also be more effective in sustaining high rates of social
interaction in their---constantly shifting, sparse and loosely
bounded---personal communities (Wellman 1979; Wellman, Carrington, and
Hall 1988; Wellman and Wortley 1990). This should also help the
individual in maintaining ties to important ``interaction foci'' (Feld
1981, 1982) where old relations are maintained and new social
connections are forged.

\emph{H4: High rates of social connectedness at a given time are
directly related to cultural taste variety in the previous time period
net of reciprocal influences of social connectedness on taste.}

\section{Data and variables}\label{data-and-variables}

The data for this study come from the \emph{Project Canada Survey}
(PCS). The PCS is a longitudinal study of the social attitudes and
cultural practices of a representative sample of Canadian citizens
extending from 1975-1995, conducted out of the University of Lethbridge.
The project began data collection in 1975 with a total 1,917
respondents. Every five years since representative samples of similar
size have been collected, with a ``core'' set of respondents that have
participated in previous surveys. Because it collected data on cultural
practices and patterns of social interaction and membership in important
foci of social interaction (i.e. voluntary association memberships) over
time for the same individuals, the PCS data provide an excellent
opportunity to deal with issues of casual order in the relationship
between cultural taste loss and life-course events associated with
network turnover.

In the following study, I use the two-wave panel covering the years 1985
and 1990. While there is data available for these same respondents for
the year 1995, I restrict myself to the core set of respondents who
participated in both of these waves (N=560).\footnote{This subset of the
  sample while not representative of the Canadian population can be used
  to monitor change in habits and practices and the determinants of
  these changes as it is attempted in the present study. Response rates
  for 1985 and 1995 where about 60\%. About 65\% of those interviewed in
  1985 where available to be re-interviewed in 1995.} . I do this for
two reasons. First, it is well know that cultural participation declines
with age (Smith 1995), so it is advantageous to measure cultural taste
with behavioral data before cultural activity begins to tail off.
Second, these two waves used a comparable coding on the taste indicators
(a slightly different coding was used in 1995) allowing an operationally
valid measure of over-time taste loss.

\emph{Data limitations}. While this data provide one of the few existing
quantitative sources of information of cultural taste over-time for a
large sample of respondents, there are some limitations of the data at
hand that deserve mention: first, there are no direct measures of
network connectivity (i.e. size of the personal network). The two
measures available, frequency of interaction with friends, and number of
associational memberships are important components of being sociable but
certainly do not represent the notion of social connectivity in its
entirety. Second, there are slight changes in the coding of the cultural
activity and friendship interaction variables across waves, with the
introduction of a slightly more elaborate scheme in 1990 (including a
``monthly'' category) and changing to never category to a composite
``hardly ever/never.'' This may compromise inferences made in
investigating taste-loss, by leading to an over-estimation of taste loss
events for instance.

I attempted to deal with this last issue by counting a taste loss event
as one that involved a shift from one of the extreme categories to the
other (the coding of which did not change across waves) and by not using
items that showed a clear effect of the shift in wording. For instance,
it is clear that for the item related to ``going to the movies''
something like this happened, with 15\% of respondents saying that they
``never'' go to the movies in 1985, but with 84\% saying that they
either ``hardly ever or never'' go to the movies in 1990. A
cross-classification of the responses to the same question across time,
reveals that this dramatic shift is mostly due to the fact that most of
the respondents who answer ``seldom'' in 1985 (55\%), where
overwhelmingly more likely to answer ``hardly ever or never'' in 1990.
The same artificial taste shift occurs for other cultural taste
indicators associated with \emph{attending events outside the home}
(going to a sports event, going to see a dramatic performance in a
theater). While some cultural practices were affected by this wording change, I nonetheless examine the full set of ten available cultural activities to provide a comprehensive test of taste stability, adjusting for this limitation analytically.

\subsection{\texorpdfstring{Variables }{Variables }}\label{variables}

\subsubsection{Cultural Taste
indicators}\label{cultural-taste-indicators}

\\emph{Cultural taste loss indicators}. The PCS collected information for
the years 1985 and 1990 related to the frequency of consumption of a
variety of leisure culture activities and mass media outlets. To
investigate the hypotheses related to cultural taste loss, I examine 
ten activities covering consumption of culture, the arts, and mass media: 
(1) listening to music, (2) going to the movies, (3) attending a play, 
(4) going to a sports event, (5) watching videos, (6) engaging in a hobby, 
(7) reading magazines, (8) following a sports team, (9) reading the newspaper, 
and (10) reading books of fiction. Respondents were asked to report
the frequency with which they engaged in each of these activities using
the following prompt:\footnote{The various codebooks for all waves of
  the PCS survey (1975, 1980, 1985, 1990 and 1995) are available online
  at: http://www.thearda.com/Archive/IntSurveys.asp.}

\emph{Thinking now of your life outside of formal work, how often would
you estimate that you: {[}i.e. listen to music{]}?}

Four categories of responses were allowed: (1) very often, (2)
sometimes, (3) Seldom and (4) Never.\footnote{As noted above, the 1990
  coded the same variables using a slightly more elaborate
  categorization. I recoded those variables to match the 1985 scheme.}
Thus for each respondent (\emph{i}), for a given cultural activity
(\emph{k}) at time (\emph{t}) is defined as their score in the
respective ordered categorical variable y\textsubscript{\emph{ikt}.}

\emph{Cultural taste breadth indicators}. To assess the hypotheses
connected to the effects of cultural on indicators of sociability and
social connectivity, I created a cultural taste breadth variable for the
1980, 1985 and 1990 waves. This was a simple additive index including
participation in a wider range of cultural activities than those used in
the cultural taste loss analysis. In addition to the four activities
discussed above, the respondent receives one point in the cultural taste
breadth scale if he or she reports engaging in the following cultural
practices either ``very often'' or ``sometimes'': (1) listening to
music, (2) going to the movies, (3) going to a theatric performance, (4)
going to a sports event outside the home, (5) watching videos at home,
(6) engaging in a hobby, (7) reading magazines, (8) following a sports
team, (9) reading the newspaper, and (10) reading a novel or book.

\subsubsection{Network turnover
indicators}\label{network-turnover-indicators}

As noted above, network theory and research provide several guidelines
as to what life-course events are associated with radical changes in the
personal network. The PCS contains two measures that tap more or less
direct transformations of the personal network: (1) longitudinal
measures of frequency of interaction with friends and (2) over-time
measures of memberships in a variety of voluntary associations
(McPherson 1983; McPherson and Smith-Lovin 1987). I also use three other
indicators of personal network turnover: (3) changes in educational
status (Suitor and Keeton 1997) (4) changes in work status (Fischer and
Oliker 1983), (5) changes in geographical location (Degenne and Lebeaux
2005), (6) changes in marital status (Fischer and Oliker 1983) and (7)
changes child rearing status (Munch, McPherson, and Smith-Lovin 1997).
Since a non-negligible part of the individual's personal network comes
from contacts that are often seen at work (Brass 1985; Ibarra 1992;
Straits 1996), we should expect that either entrance into or exit from
the labor force, or shifts from part-time to full-time work, should have
a clear impact on the size and composition of the personal network over
time (Bidart and Lavenu 2005). The same can be said of geographical
mobility (Relish 1997), with geographically mobile respondents more
likely to be subject to radical transformation of their personal
networks from those who are not mobile.\footnote{The nominal employment
  status from which the binary indicator for changes in work status is
  constructed from has six categories (1) full time, (2) part time, (3)
  unemployed, (4) retired, (5) keeping house, and (6) other. The
  geographical mobility item is based on the region in which respondent
  reported residing in each wave. The marital status variable has five
  categories: (1) married, (2) cohabiting, (3) divorced, (4) widowed and
  (5) never married.} Finally, a large body of research shows that
family status transitions, such as getting married or getting divorced
(Gerstel 1988; Leslie and Grady 1985; Milardo 1987; Pahl and Pevalin
2005) and experiencing widowhood (Anderson 1984; Morgan, Carder, and
Neal 1997; Morgan, Neal, and Carde 1997), have significant effects on
the social networks of men and women. We should thus expect that
transformations in family life, such as getting divorced, getting
married or having children (Belsky and Rovine 1984; Bott 1971; Munch,
McPherson, and Smith-Lovin 1997), should result in direct effects on the
size and composition of the personal network.

\subsubsection{Social connectedness
indicators}\label{social-connectedness-indicators}

In order to test the cultural capital hypotheses (3 and 4), I use two
indicators of social connectivity in the most recent wave of the survey:
(a) rates of social interaction with friends in 1995 and (b) the ability
to remain connected to important foci for social interaction from one
period to the next, such as voluntary associations. The social
interaction measure was obtained from the response to the following
question:

\emph{Thinking now of your life outside of formal work, how often would
you estimate that you - Spend time with friends?}

In 1985, the allowed responses were: (1) very often, (2) sometimes, (3)
Seldom and (4) Never. In 1990, the respondents chose from five response
categories: (1) very often, (2) sometimes, (3) Seldom (4) Monthly, (5)
Hardly ever or never. Because the models used below are designed to
predict interval variables, and due to the fact that there is no reason
to expect that the numerical categories used in the survey clearly
correspond to the naïve numerical ordering (or that in the case of the
categories used in 1990, that the categories are in the \emph{correct}
order; is ``monthly'' more or less that ``sometimes'') I use association
models for cross-classifications (Goodman 1979) to induce the ``latent
ordering'' of the cross-classification of responses to the frequency of
interaction with friends across the two waves (Clogg 1982). I use these
scores and not the raw categories as the dependent variables in the
analyses reported below. The results are shown in Table 1. The table
shows that in both 1990 and 1985, the responses do align themselves into
a unidimensional continuum (the association model fits the data well;
\emph{G\textsuperscript{2}}=1.29, \emph{p}\textless0.73), although in
1990, the distance between very often and sometimes is not as large as
that between sometimes and seldom.

Retention of connections to important foci interaction is measured using
a scale that simply counts the types of voluntary associations in each
wave of the survey: these include private clubs, sport clubs, service
organizations and hobby clubs. These types of memberships are the kinds
of interaction foci in which mass media and arts related culture
mediated social interaction is most likely to play a key role in the
formation and maintenance of social contacts (Long 2003), and thus on
the likelihood that the individual will retain or drop the membership
(McPherson, Popielarz, and Drobnic 1992; Popielarz and McPherson 1995).

\includegraphics[width=6.01319in,height=3.10417in]{media/image1.emf}

Figure 1

\subsubsection{Sociodemographic
controls}\label{sociodemographic-controls}

The models shown below include the following sociodemographic controls:
gender, age, education, city size, employment status, and marital
status, presence of children in the household and city size. The coding
and descriptive statistics of all of the variables mentioned above are
shown in Table 4 of the Appendix.

\subsection{\texorpdfstring{Network driven taste loss
}{Network driven taste loss }}\label{network-driven-taste-loss}

\subsubsection{Analytic Strategy}\label{analytic-strategy}

The cultural taste loss model can be thought of as a discrete time
analogue to the usual continuous time event history model for two time
periods (Allison 1982), with the dependent variable equaling 1, if a
cultural taste loss event is observed and zero otherwise. The
coefficients of interest can then be estimated using the following
logistic regression equation:

\(\ln\) (1)

Where \emph{p\textsubscript{ik}} is the probability that individual
(\emph{i}) reports that he or she ``never'' or ``seldom'' engages in
cultural activity (\emph{k}) in 1990 after reporting that he or she did
so ``very often'' in 1985 and \emph{M} is the number of sociodemographic
factors included in the equations. The dependent variable is thus our
measure that there has been a shift in taste for cultural form \emph{k}
in the 5 year period. \emph{Z} in (1) is an \emph{i} x \emph{M} design
matrix of network turnover indicators. The matrix is constructed from
each of the sociodemographic and positional variables
(\emph{w\textsubscript{im}}) whose over-time fluctuation (i.e. going
from employed in the first time period to retired in the second) is
taken to be likely to lead to cultural taste loss The Z matrix is coded
according to the following scheme:

\(z_{im}^{1985} = 0\) (2)

\(z_{im}^{1990} = 0\ \ \ \overset{}{\leftrightarrow}\ w_{im}^{1990} = \ w_{im}^{1985}\)
(3)

\(z_{im}^{1990} = 1\ \ \ \overset{}{\leftrightarrow}w_{im}^{1990} \neq w_{im}^{1985}\)
(4)

Thus for each entry \emph{Z\textsubscript{m}}, of the network turnover
matrix, a 1 is recorded if a change in sociodemographic status \emph{w}
is observed in 1990 and a zero is recorded otherwise.\footnote{The
  voluntary association membership change variable records a one if the
  individual reports a change in membership status (taken from a binary
  variable that equals one if the respondent reports belonging to the
  organization and zero otherwise) for any one of eight types of
  voluntary associations. Therefore this is a dichotomized version of an
  originally categorical count variable. The maximum observed number of
  memberships changes was 6 and the minimum was zero, with a mean of 1.7
  and a standard deviation of 1.3.} In the regression model, \emph{γ} is
thus an \emph{M} x 1 vector of coefficients associated with the effect
of network turnover for \emph{m}\textsuperscript{th} covariate on
cultural taste loss. If the network hypothesis is correct (hypothesis 2
above), then we should observe \emph{γ}\textgreater0 for each of the
positional shift indicators. Finally, the \emph{X\textsubscript{ik}} is
equation (1) is a matrix of time invariant (i.e. gender and ethnicity)
and time variant (i.e. education and age) individual sociodemographic
characteristics and \emph{β} is a \emph{k} x 1 vector of coefficients
corresponding to the effects of these variables on the odds of cultural
taste loss.\footnote{Cultural taste and sociodemographic data are
  available for 560 individuals in the two time periods. Since each
  respondent yields a case in each of the two panels, the maximum
  corresponding number of observations for the following analysis is
  (560 x 2=1120). After deleting observations with missing values we are
  left with N=1116 in the discrete-time logistic regression models with cluster-robust standard errors.}

\subsubsection{Results}\label{results}

\emph{Frequency of Taste Loss.} \textbf{Error! Reference source not
found.} shows the proportion of respondents who experienced a ``taste
loss'' event from 1985 to 1995. Consistent with the cultural capital
hypothesis, across all four cultural taste domains, taste loss is a
relatively rare event. This is the case whether this is defined in a
``strong'' (front ``very often'' in the first wave to ``never'' in the
second) or in a weaker form (from ``very often'' to \emph{either}
``seldom'' or ``never''). Less than 1\% of respondents go from having a
strong taste in listening to music or reading the newspaper to losing
that taste in the five year period. A relatively larger number of
respondents lose the taste for reading books (1.07\%) or following
sports (1.43\%) but the main empirical finding here is that the great
majority of respondents retain their tastes (≈98\%). When taste loss is
defined in a weaker form the worst retention rate drops to about 95\%
for reading books. This supports the cultural capital assumption that
tastes are relatively stable over time, a rate of stable that appears to
exceed that which can be observed for personal networks (Suitor,
Wellman, and Morgan 1997; Wellman \emph{et al.} 1997).

\emph{Taste loss Models.} The coefficient estimates of the discrete-time
event history models are shown in Table 2. All of the models have as the
dependent variable a taste loss event defined according to the ``weak''
criterion above. Consistent with the expectations derived from the
network-based taste transmission theory (Hypothesis 2), I find
that---however rare---the cultural taste loss that can be observed is
largely due to positional and social status changes associated with
network turnover. Thus, the best predictors of cultural taste loss are
changes in the frequency of \emph{social interaction with friends,}
changes in the \emph{number of organizational memberships}
(\emph{γ\textsubscript{k }}\textgreater0, \emph{p}\textless0.05 for
\emph{all} cultural activities). Other important determinants of taste
loss are status transitions that have been shown in the literature to be
clearly associated with network turnover such as changes in educational
standing (\emph{p}\textless0.05 for all cultural activities except
reading books and reading magazines) and changes in geographic location
(\emph{γ} also statistically significantly different from zero for all
four activities). The strongest effects in terms of magnitude, is in
accord with the view that tastes change when the personal network is
transformed that which pertains to changes in social interaction.

Changes in work status also have the expected positive effects on taste
loss, although they are only statistically significant for three out of
the four cultural activities (reading the paper is not affected by this
factor). The effect of work status on taste for sports in consistent
with Erickson's (1996) observation of the key role of ``sports culture''
in smoothing interaction across class lines in the workplace. The
weakest effect of the entire set of network turnover indicators on taste
loss belongs pertains to those associated with changes in family status.
Changes in marital status is not a useful predictor of taste loss for
three out of the four cultural activities, once other sources of network
turnover are held constant, although it still predicts changes in taste
for listening to music. Changes in the number of children in the
household on the other hand do not have a discernible effect on cultural
taste loss. Possible reasons for the paucity of effects of changes in
marital status include measurement error due to underestimating family
volatility (respondents who report being married in both waves are
counted as not having experienced a change, even though they might be
married to different people or may have experienced other changes in the
interim), or the fact that there is little net effect of family status
changes left after the primary relational outcomes of this changes
(shift in the frequency of interaction with friends) are held constant.

In all looking at the relatively high χ\textsuperscript{2} value for
overall model fit and the pseudo-R-squared values, it is safe to say
that the network hypothesis is strongly confirmed by these data. While
the \emph{rates} of cultural taste loss are not consistent with standard
network theory, the \emph{sources} of taste loss when it does happen are
largely certainly those that we would expect to be responsible for these
changes: transformation of the personal network due to life-course
changes associated with network turnover.

\subsection{Longitudinal Culture conversion
model}\label{longitudinal-culture-conversion-model}

\subsubsection{Analytic Strategy}\label{analytic-strategy-1}

In this section I estimate the possible casual effect of cultural
capital (as measured by taste diversity) on the maintenance of variables
associated with network connectivity over time. There are particular
theoretical and empirical challenges in estimating the effects of
cultural taste on indicators of social connectivity and frequency of
interaction. First, there is the issue of \emph{unobserved individual
heterogeneity}: any effect of cultural taste on social interaction
observed across panels might be driven by an unobserved, time-invariant
individual-level factor that is correlated with both variables. Second
there is the issue of \emph{simultaneity bias}: if cultural taste
breadth is affected by frequency of social interaction with friends and
connectivity to important relational foci, then any estimate of the
reciprocal effect of taste on those variables will be biased and
inconsistent.

I deal with both of these issues in turn by presenting three different
specifications, all of which take advantage of the availability of panel
data at two points in time:

\emph{Lagged independent variables model}. First, I present I lagged
independent variable model, in which the effects of cultural taste
breadth 1985 on measures of social network connectivity in 1990 are
estimated holding constant sociodemographic factors measured in 1985.
While this strategy takes explicit advantage of the longitudinal aspect
of the data, this specification is open to bias produced by any
unmeasured, time-invariant individual level factor not included in the
model that is constant over time and which thus spuriously produces the
cross-temporal association of cultural taste breadth and network
connectivity. The longitudinal nature of the data allows us to deal
effectively with the issue of unobserved individual heterogeneity by
modeling it directly. In the following I opt for a random-effects
modeling framework. This model is:

\({E(Y}_{i}^{1990}) = a + \gamma C_{i}^{1985} + \sum_{}^{}{\beta x_{i}^{1985} + e_{i}}\)
(5)

Where \emph{Y\textsubscript{i}} is an indicator of social connectivity
measured in 1990, \emph{C\textsubscript{i}} is a measure of cultural
taste breadth measured in 1985 and x is a vector of sociodemographic
control variables measured in 1985.

\emph{Random-effects models}. I also present a model that accounts for
unobserved heterogeneity by including an individual level random effect.
However, while the random-effects model takes care of the unobserved
heterogeneity issue, thus insuring that any observed effect of cultural
taste breadth on social connectivity is not spuriously caused by some
underlying individual-level factor that is constant across time, it is
still open to the simultaneity bias caused by reciprocal causation going
from social connectivity to taste breadth (Lizardo 2006). This model can
be written as:

\({E(Y}_{it}^{}) = a + \gamma C_{it}^{} + \sum_{}^{}{\beta x_{it}^{} + {u_{i} + e}_{it}}\)
(6)

Where everything is as in (5), except that now the observation pertains
to the individual-time period, as shown by the inclusion of the \emph{t}
subscript. This means that each individual yields two observations and
the random-effect (\emph{u\textsubscript{i}}) is parameterized as being
constant across time but varying across individuals. Finally,
\emph{e\textsubscript{it}} is a classical random disturbance that varies
across time and across individuals.

\\emph{Within-Between (Mundlak) Poisson Models}. To take this issue 
into account, I use a within-between mixed-effects Poisson model (often referred 
to as a Mundlak model) for the panel data. This model accounts for both 
time-invariant unobserved heterogeneity (like a fixed-effects model) while still 
allowing the estimation of group-level differences, avoiding the strict assumptions 
of standard random-effects models. Because the dependent variables (number of 
friends and organizational memberships) are count variables, a Poisson specification 
is appropriate.\footnote{This model is the same as (6) except that the
  measure of cultural taste breadth comes from the \emph{predicted
  value} of a regression equation that includes all of the control and
  independent variables and the value of cultural taste breadth in 1980.}

\subsubsection{Results}\label{results-1}

The culture conversion model (Lizardo 2006) implies that individuals who
have command of a wide variety of tastes will be able to sustain larger
and sparser personal networks over time, allowing them to reap the
benefits of those social connections in the form of higher than average
rates of social interaction and a higher ability to remain connected to
important social foci where new contacts are made and old ones are
sustained. The results of the lagged, random effects and instrumental
variables regression models of the effect of cultural capital (in the
form of ``omnivore'' taste) in the first time period on network outcomes
in the second time period are shown in Table 3.

The results shown in the first model of the table are consistent with
the culture conversion model: I find that persons who display a wide
variety of tastes in the initial time period sustain higher frequencies
of social interaction five years later \emph{even after holding constant
individual level predictors of interaction frequency in the first time
period} (γ\textgreater0, \emph{p}\textless0.05). While not definitive
evidence of ``causality'' these results serve to cement in a more
straightforward longitudinal framework those reported by Lizardo (2006)
using two-stage methods on cross-sectional data, thus serving to more
firmly establish the primary empirical implication of the culture
conversion model: broad cultural tastes are used to sustain larger and
in this case more active patterns of interpersonal engagement over time.
Furthermore, we can see in the fourth column of the table having a wide
repertoire of taste is not only useful in sustaining higher levels
informal social connectivity to other members of the personal community
network over time, but that it is also helps in raising the chances of
remaining tied to important social foci of interaction such as voluntary
associations (\emph{γ}\textgreater0, \emph{p}\textless0.05). As Feld
(1981) has noted these interaction foci are important since they serve
as the primary social mechanism that allows the individual to form
social connections which serve to integrate her into the larger social
system. Because recruitment (and the over-time stability of membership)
into this foci is largely a matter of keeping alive the network ties
that bind the individual to other members of the group (Popielarz and
McPherson 1995), it is no surprise that whatever resource helps the
individual to sustain social connections to other persons will also
appear to sustain his or her connection to these larger social entities.

Are the results from the lagged models the spurious by-product of
unmeasured individual heterogeneity or simultaneity bias? The models
shown in the second and fifth columns of the table show the results
after taking into account unobserved individual heterogeneity by
modeling it within a random effects framework. The basic result stands:
persons who engage in a wide variety of cultural activities are better
able to sustain higher levels of social connectivity both in terms of
interaction with members of the ``personal community'' (Wellman 1979,
Wellman et al 1988) and in terms of connection with organizational
interaction foci. Furthermore, as shown in the third and sixth columns
of the table, the results remain the same (in fact become stronger,
suggesting that the simultaneity bias was depressing the magnitude of
the cultural taste coefficient downwards) after we deal with both the
simultaneity and unobserved heterogeneity bias by using cultural taste
breadth in 1980 as a Within-Between model framework. Thus we
can conclude that in the very same way that being a cultural
``omnivore'' helps to sustain social connections with friends, cultural
variety helps the individual keep their connections to voluntary
association foci and sustain high levels of informal sociability with
members of the personal community network.

\subsection{Discussion and conclusion}\label{discussion-and-conclusion}

This paper advances current theory and research in several ways. First,
it provides for the first time conclusive empirical evidence that
pertains directly to the empirical adequacy of different assumptions
regarding the question of the stability of cultural tastes over time.
The question of taste stability is important, since it bears on the
direct empirical implications of the two most dominant sociological
theoretical frameworks that help us explain the relationship between
social ties and cultural resources: network theory whether in its
relational (Erickson 1996) or ecological (Mark 1998a, 199b; McPherson
2004) variants, and cultural capital theory (Bourdieu 1984, 1986).
Second, this paper produces direct empirical evidence of another crucial
assumption of network-theoretic views of taste formation, transmission
and decay: that of the effect of relational disruption on taste loss. I
tackle this question by attempting to link longitudinal shifts in
patterns of social interaction, connection to important social foci
(Feld 1981) and life-course transitions associated with network turnover
to changes in cultural taste. This represents the first time that this
most crucial of implicit assumptions of network theories of taste (Mark
1998b) receives empirical treatment in a research design where issues of
casual direction can be dealt with satisfactorily. Finally, this paper
builds on previous research on the role of cultural capital and other
types of cultural resources on the creation and maintenance of social
ties (DiMaggio 1987, 1993; DiMaggio and Mohr 1985; Lizardo 2006; Mische
2003) by providing relatively unambiguous evidence of the reciprocal
casual effect of durable patterns of cultural taste on the maintenance
of patterns of social interaction and connection to social foci
over-time.

The results show that the cultural taste instability corollary of
traditional network accounts appears not to be operative in this data.
Instead, taste for music, newspaper reading, books, and sports show high
degrees of stability, a result consistent with the assumptions of a
cultural-capital model. However, even though the taste-stability
corollary is not supported by these data, the network theoretic
assumption that cultural taste loss is more likely to be due to events
associated with the disruption of personal networks receives very strong
support. Across all four cultural forms included in the analysis taste
loss events are strongly associated with changes in interaction
frequency with friends, alteration in patterns of connection to
important social foci, and disruptions associated with geographic
mobility and changes in work status. Finally, a more stringent test of
the cultural conversion hypothesis than that offered by Lizardo (2006)
produces robust substantiation for the theory. Even after taking into
account both simultaneity and unobserved heterogeneity bias, we find
that individuals who command a wide variety of cultural tastes are more
likely to sustain higher levels of social interaction with friends and
higher levels of connectivity to important interaction settings (such as
voluntary association memberships) through the five year period. This
effect it is important to underscore appears not to be due to the
reciprocal effect of connectivity on cultural taste breadth.

The results presented here have significant implications for how we
theorize the statics and dynamics of cultural taste formation,
transmission and decay. Rather than the generic idea found in the
traditional network model of the individual who is influenced by the
tastes of those persons that he or she is connected to, the cultural
capital derived hypotheses of relative taste stability and use of
cultural taste for the formation of new relationships point to two
important additional considerations usually ignored in traditional
network-based accounts.

\emph{Choice-homophily of network contacts based on pre-existing
cultural dispositions}. Standard network theories of taste transmission
draw on the notion of ``homophily'' (i.e. the empirical generalization
that close friends tend to be similar in sociodemographic
characteristics) to explain the fact that cultural tastes and practices
tend to cluster in certain areas of social space (Mark 2003). From the
standard point of view homophilious social connections based on
socio-demographic similarity \emph{pre-exist} individual patterns of
taste and this explains why cultural practices ``lump'' in certain areas
of social space. That is, the logic of network-based theories leads to
them to attempt to explain the fact that cultural practices occupy
delimited niches in social space by pointing to two interrelated
processes: (1) tastes tend to be associated with social background and
(2) individuals tend to associate with those who are
socio-demographically similar.

Yet, there is sufficient qualitative evidence (Cardon and Granjon 2005;
Long 2003; Wali, Severson, and Longoni 2002)---and as this paper shows,
quantitative evidence (see also Lizardo 2006)---to suggest that just in
the very same that tastes are transmitted through pre-existing network
ties, \emph{new} social ties are formed \emph{on the basis of
pre-existing cultural dispositions} (Lazarsfeld and Merton
1954).\footnote{See Illouz 1998 on romantic partnerships and cultural
  capital DiMaggio 1996 and Ostrower 1998 on taste for the high arts and
  social relations among elites and Bennett 2004 on taste for popular
  music and social network formation in online communities.} This is the
obverse of the traditional network-theoretic emphasis on the formation
of new cultural tastes on the basis of preexisting social ties. In fact,
given the relative stability of cultural tastes in comparison to network
connections, it is likely that this last process may in fact occur more
often that the former process (Lizardo 2006). In this respect (and in
contradiction to standard network accounts), the ``lumpiness'' of
cultural tastes in social space may be actually instead be the result of
an active---but not necessarily conscious---process of contact selection
and relationship fine-tuning \emph{keyed around cultural compatibility}
(Lazarsfeld and Merton 1954)\emph{.}

This proposition is concordant with Burt's (2000) observation of the
``liability of newness'' of recently formed ties, and with recent
studies of changes in social networks of young adolescents and young
adults (Bidart and Degenne 2005; Bidart and Lavenu 2005; Cardon and
Granjon 2005). Notice that if this is correct, then sociodemographic
homophily could in the majority of cases be a \emph{by-product} of
cultural compatibility and not the other way around (this cultural
compatibility may of course be itself due to socio-demographic of the
family of origin {[}Bourdieu 1984; Illouz 1998{]}). This is an insight
that was always implicit in the theoretical logic of dynamic simulations
of network formation based on shared cultural knowledge (Carley 1991),
but which has been occluded in recent theory and research due to the one
sided emphasis put on the socio-structural determination of shared
tastes.

\emph{Class-based heterogeneity in contact choice and structural
determination of tastes}. Another insight that is brought about by the
current research is the fact that the extent to which contacts are
subject to a process of active selection based on taste itself varies
according to the individual's command of cultural, educational and
socio-economic resources. Put simply, high cultural capital individuals,
in particular those employed in prestigious occupations that demand of
extensive degrees of cognitive and cultural skill not only have larger
and sparser social networks (Lizardo 2006), but the proportion of their
networks that is composed of constantly changing, loosely bounded social
ties is larger and more heterogeneous than that of culturally
disadvantaged individuals employed in routine, low-complexity
occupations (Coser 1991; Fischer 1982). In this respect, the endogenous
effect of cultural capital in both the ability to select among different
social ties and the ability to \emph{reach} across the social structure
to form those ties, is itself partially determined by position in social
space.

This is important, because the dominant way that ``social space'' has
been conceptualized by network and organizational theorists (i.e.
Hannan, Carroll, and Pólos 2003; McPherson 2004) in relation to culture
has been---following Blau (1977)---as a multidimensional manifold
composed of variables such as age and occupational prestige. The
sociodemographic niche for a particular cultural or social form can then
be given a metric definition in terms of these variables. Social
distance among two points in these space is conceived as being
proportional to sociometric distance in social networks, that as being a
proxy for the probability that two individuals will interact in real
life (Mark 1998a). This presumes that the socio-demographic covariates
that define social space are strictly exogenous (i.e. they do not
interact in the statistical sense) with the propensity to deploy
cultural knowledge reaped from consuming the cultural form in question
in interaction to form new social relationships with persons located at
relatively distant locations in social space. Traditional network models
also assume that there is no interaction between the occupancy of
privileged locations in social space and the relative propensity to
acquire or learn about specific cultural forms from proximate social
contacts. This process is seen as being a generic dynamic applicable to
all individuals regardless of their socio-demographic profile.

From the point of view of the culture-conversion model on the other hand
the axes of social space as defined above are clearly composed of
variables some of which (i.e. educational attainment) themselves are
implicated in the extent to which tastes may be driven by social network
dynamics. Thus, from this point of view, occupying a privileged position
in social space plays an important role in the degree to which
individuals are able to loosen the limiting effects of network-based
social constraint. This happens due the greater ability of those who
posses higher levels of educational and socio-economic resources to
deploy a more heterogeneous repertoire of cultural knowledge in everyday
interaction. This allows them to form and sustain social relationships
that cut-across social dimensions of stratification and differentiation
(DiMaggio 1987), as well permitting them to create and maintain intimate
ties with those who share their knowledge and tastes. For this reason,
the dimensions of ``social space'' cannot be considered strictly
exogenous to the cultural contents whose sociodemographic niche they
define. Instead differential command of these cultural contents on the
part of individuals may themselves ``warp'' these dimensions, such that
individuals located at some privileged locations in this manifold find
it easier to traverse a large distance in that space and thus connect to
a wide range of others (Erickson 1996; Lin and Dumin 1986). Individuals
located in less privileged areas of the space, on the other hand, are
socially and culturally distant from most individuals and cannot connect
very well across their locally bounded social circles (Bryson 1997).

There is plenty evidence that is consistent with these considerations.
For instance Wellman (1979; Wellman et al 1988; Wellman and Wortley
1990) and Fischer (1982) show that as socio-economic and educational
status increases, individuals are much more likely to have
role-segmented, specialized relationships in which a clear
differentiation is made between those ties designed for sociability and
companionship (usually labeled as ``friends'') and social ties that are
the result of non-negotiable statuses such as kinship ties. These lines
of research on personal community networks also show that individuals
are much more likely to prefer to obtain companionship from voluntary
rather than involuntary ties. These are the social contacts that the
culture conversion model predicts are likely to be formed and sustained
through pre-existing forms of cultural compatibility. These are also the
social relationships that are likely to be sustained through interaction
rituals in which the symbolic goods coming from artworlds and cultural
industries figure as the primary subject of conversation (Cardon and
Granjon 2005; DiMaggio 1987). Finally, these are also the same social
ties that are continually undergoing high rates of turnover and being
constantly recycled and reformed, and which thus require the recurrent
deployment of cultural capital in interaction. This goes a long way
toward explaining the almost insatiable demand for novel symbolic goods
among high cultural capital class fractions (Douglas and Isherwood
1996).

In essence, for a relatively privileged cultural and economic elites,
social connections to culturally affine others have come to be
disembedded from place, kinship and ascription (DiMaggio and Mohr 1985;
Holt 1998), and have taken the form of increasingly longitudinally
unstable and thus constantly changing personal community networks
(Wellman et al 1988). This increases the chances that differential
command of the ability to consume a wide variety of cultural forms
becomes a \emph{driver} of the recurrent formation of new network
relations. This is less likely to be the case for those individuals who
are restricted to a more homogenous, kin-based, and locally constrained
set of core alters who change relatively little over time. Consistent
with this claim, Lizardo (2006) shows that after holding cultural
capital constant, the positive effect of education and socio-economic
status on social network size is reduced significantly. This has the
curious implication that traditional network theory may more accurately
describe the process of cultural taste formation among low cultural
capital individuals subject to high levels of network constraint (Burt
2005) while the process of culture conversion may be a more accurate
representation of the deployment and operation of cultural taste among
culturally advantaged elites.

\subsection{References}\label{references}

Alexander, Victoria D. 2003. \emph{Sociology of the Arts: Exploring Fine
and Popular Forms}. Malden: Blackwell.

Allison, Paul D. 1982. " Discrete-Time Methods for the Analysis of Event
Histories." \emph{Sociological Methodology} 13: 61-98.

Anderson, Trudy B. 1984. "Widowhood as a Life Transition: Its Impact on
Kinship Ties." \emph{Journal of Marriage and the Family} 46:105-114.

Belsky, Jay and Michael Rovine. 1984. "Social-Network Contact, Family
Support, and the Transition to Parenthood." \emph{Journal of Marriage
and the Family} 46:455-462.

Bennett, Andy. 2004. "New Tales from Canterbury: the Making of Virtual
Scene." Pp. 205-220 in \emph{Music Scenes: Local, Translocal and
Virtual}, edited by A. Bennett and R. A. Peterson. Nashville: Vanderbilt
University Press.

Bennett, Tony, Michael Emmison, and John Frow. 1999. \emph{Accounting
for Tastes: Australian Everyday Cultures}. Cambridge: Cambridge
University Press.

Berkowitz, S. D. 1982. \emph{An Introduction to Structural Analysis: The
Network Approach to Social Research}. Toronto: Butterworths.

\_\_\_\_\_\_. 1988. "Afterword: Toward a Formal Structural Sociology."
Pp. 477-497 in \emph{Social Structures a Network Approach}, edited by B.
Wellman and S. D. Berkowitz. Cambridge: Cambridge University Press.

Bidart, Claire and Alain Degenne. 2005. "Introduction: the dynamics of
personal networks." \emph{Social Networks} 27:283-287.

Bidart, Claire and Daniel Lavenu. 2005. "Evolutions of personal networks
and life events." \emph{Social Networks} 27:359-376.

Blau, Peter M. 1977. "A Macrostructural Theory of Social Structure."
\emph{American Journal of Sociology} 83:26-54.

Borgatti, Stephen P. 2005. "Centrality and network flow." \emph{Social
Networks} 27:55-71.

Bott, Elizabeth. 1971. \emph{Family and Social Network}. London:
Tavistock Institute of Human Relations.

Bourdieu, Pierre. 1984. \emph{Distinction: a Social Critique of the
Judgment of Taste}. Translated by R. Nice. Cambridge: Harvard University
Press.

\_\_\_\_\_\_. 1986. "The Forms of Capital." Pp. 241-258 in
\emph{Handbook of Theory and Research for the Sociology of Education},
edited by J. Richardson. New York: Greenwood Press.

\_\_\_\_\_\_. 1990. \emph{The Logic of Practice}. Cambridge: Polity
Press.

\_\_\_\_\_\_. 2000. \emph{Pascalian Meditations}. Stanford: Stanford
University Press.

Brass, Daniel J. 1985. "Men\textquotesingle s and
Women\textquotesingle s Networks: A Study of Interaction Patterns and
Influence in an Organization." \emph{The Academy of Management Journal}
28:327-343.

Bryson, Bethany. 1997. "What about the Univores? Musical Dislikes and
Group-Based Identity Construction among Americans with Low Levels of
Education." \emph{Poetics} 25:141-156.

Burt, Ronald S. 2000. "Decay Functions." \emph{Social Networks} 22:1-28.

\_\_\_\_\_\_. 2002. "Bridge Decay." \emph{Social Networks} 24: 333-363.

\_\_\_\_\_\_. 2005. \emph{Brokerage and Closure}. Oxford: Oxford
University Press.

Cardon, Dominique and Fabien Granjon. 2005. "Social networks and
cultural practices A case study of young avid screen users in France."
\emph{Social Networks}:301-315.

Carley, Kathleen M. 1991. "A Theory of Group Stability." \emph{American
Sociological Review} 56:331-354.

\_\_\_\_\_\_. 1995. "Communication Technologies and Their Effect on
Cultural Homogeneity, Consensus, and the Diffusion of New Ideas."
\emph{Sociological Perspectives} 38: 547-571.

Clogg, Clifford C. 1982. "Using Association Models in Sociological
Research: Some Examples." \emph{American Journal of Sociology}
88:114-134.

Collins, Randall. 2004. \emph{Interaction Ritual Chains}. Princeton:
Princeton University Press.

Coser, Rose Laub. 1991. \emph{In Defense of Modernity: Role Complexity
and Individual Autonomy}. Stanford: Stanford University Press.

Crane, Diana. 1993. "High Culture versus Popular Culture Revisited:
Toward a Reconceptualization of Recorded Cultures." Pp. 58-74 in
\emph{Cultivating Differences: Symbolic Boundaries and the Making of
Inequality}, edited by M. Lamont and M. Fournier. Chicago: Chicago
University Press.

Crook, Christopher J. 1997. "The Dimensionality of
Stratification-Related Cultural Practices in Australia." \emph{Journal
of Sociology} 33:226-238.

Degenne, Alain and Marie-Odile Lebeaux. 2005. "The dynamics of personal
networks at the time of entry into adult life." \emph{Social Networks}
27:337-358.

DiMaggio, Paul. 1987. "Classification in Art." \emph{American
Sociological Review} 52:440-455.

\_\_\_\_\_\_. 1993. "Nadel\textquotesingle s Paradox Revisited:
Relational and Cultural Aspects of Organizational Structure." Pp.
118-142 in \emph{Networks and Organizations}, edited by N. Nohria and R.
G. Eccles. Boston: Harvard Business School Press.

\_\_\_\_\_\_. 1996. "Are art-museum visitors different from other
people? The relationship between attendance and social and political
attitudes in the United States." \emph{Poetics} 24:161-180.

\_\_\_\_\_\_. 2000. "Social Structure Institutions and Cultural Goods:
The Case of the United States." Pp. 38-62 in \emph{The Politics of
Culture}, edited by G. Bradford, M. Gary, and G. Wallach. New York: New
Press.

DiMaggio, Paul and John Mohr. 1985. "Cultural Capital, Educational
Attainment, and Marital Selection." \emph{American Journal of Sociology}
90:1231-1261.

Dimaggio, Paul and Michael Useem. 1978. "Social Class and Arts
Consumption: The Origins and Consequences of Class Differences in
Exposure to the Arts in America." \emph{Theory and Society} 5:141-161.

Douglas, Mary and Baron Isherwood. 1996. \emph{The World of Goods:
Towards an Anthropology of Consumption}. New York: Routledge.

Erickson, Bonnie H. 1982. "Networks, Ideologies and Belief Systems." Pp.
159-172 in \emph{Social Structures and Network Analysis}, edited by P.
Marsden and N. Lin. Beverly Hills, CA: Sage.

\_\_\_\_\_\_. 1988. "The Relational Basis of Attitudes." Pp. 99-121 in
\emph{Social Structures a Network Approach}, edited by B. Wellman and S.
D. Berkowitz. Cambridge: Cambridge University Press.

\_\_\_\_\_\_. 1996. "Culture, Class, and Connections." \emph{American
Journal of Sociology} 102:217-252.

Feld, Scott L. 1981. "The Focused Organization of Social Ties."
\emph{American Journal of Sociology} 86:1015-1035.

\_\_\_\_\_\_. 1982. "Social Structural Determinants of Similarity among
Associates." \emph{American Sociological Review} 47:797-801.

Fine, Gary Alan. 1979. "Small Groups and Culture Creation: The
Idioculture of Little League Baseball Teams." \emph{American
Sociological Review} 44: 733-745.

Fischer, Claude S. 1982. \emph{To Dwell among Friends: Personal Networks
in Town and City}. Chicago: University of Chicago Press.

Fischer, Claude S. and Stacey J. Oliker. 1983. "A Research Note on
Friendship, Gender, and the Life Cycle." \emph{Social Forces}
62:124-133.

Frith, Simon. 1998. \emph{Performing Rites: On the Value of Popular
Music}. Cambridge: Harvard University Press.

Gerstel, Naomi. 1988. "Divorce and Kin Ties: The Importance of Gender."
\emph{Journal of Marriage and the Family} 50:209-219.

Goodman, Leo A. 1979. "Simple Models for the Analysis of Association in
Cross-Classifications having Ordered Categories." \emph{Journal of the
American Statistical Association} 74:537-552.

Granovetter, Mark. 2002. " theoretical agenda for economic sociology."
Pp. 35-60 in \emph{The new economic sociology: Developments in an
emerging field}, edited by M. F. Guillén, R. Collins, P. England, and M.
Meyer. New York: Russell Sage.

Hannan, Michael T., Glenn R. Carroll, and László Pólos. 2003. "The
Organizational Niche." \emph{Sociological Theory} 21:309-340.

Holbrook, Morris B. and Robert M. Schindler. 1989. "Some Exploratory
Findings on the Development of Musical Tastes." \emph{Journal of
Consumer Research} 16:119-124.

Holt, Douglas B. 1998. "Does Cultural Capital Structure American
Consumption?" \emph{Journal of Consumer Research} 25:1-25.

Holt, Douglas B. 1997. "Distinction in America? Recovering Bourdieu's
Theory of Tastes from Its Critics." \emph{Poetics} 25:93-120.

Homans, George C. 1950. \emph{The Human Group}. New York: Harcourt
Brace.

Ibarra, Herminia. 1992. "Sex Differences in Network Structure and Access
in an Advertising Firm." \emph{Administrative Science Quarterly}
37:422-447.

Illouz, Eva. 1998. \emph{Consuming the Romantic Utopia: Love and the
Cultural Contradictions of Capitalism}. Berkeley, CA: University of
California Press.

Kane, Danielle. 2004. "A network approach to the puzzle of
women\textquotesingle s cultural participation." \emph{Poetics}
32:105-127.

Kracman, Kimberly. 1996. "The Effects of School-Based Arts Instruction
on Attendance at Museums and the Performing Arts." \emph{Poetics}
24:203-218.

Lang, Frieder R. 2000. "Endings and Continuity of Social Relationships:
Maximizing Intrinsic Benefits within Personal Networks when Feeling Near
to Death." \emph{Journal of Social and Personal Relationship}
17:155-182.

Lazarsfeld, Paul F. and Robert K. Merton. 1954. "Friendship as a social
process: A substantive and methodological analysis." Pp. 18-66 in
\emph{Freedom and Control in Modern Society}, edited by M. Berger, T.
Abel, and C. Page. New York: Van Nostrand.

Leslie, Leigh A. and Katherine Grady. 1985. "Changes in
Mothers\textquotesingle{} Social Networks and Social Support following
Divorce." \emph{Journal of Marriage and the Family} 47:663-673.

Lin, Nan and Mary Dumin. 1986. "Access to Occupations through Social
Ties." \emph{Social Networks} 8:365--385.

Lizardo, Omar. 2006. "How Cultural Tastes Shape Personal Networks."
\emph{American Sociological Review} 71:778-807.

Long, Elizabeth. 2003. \emph{Book Clubs: Women and the Uses of Reading
in Everyday Life}. Chicago: University of Chicago Press.

Mark, Noah P. 1998a. "Beyond Individual Differences: Social
Differentiation from First Principles." \emph{American Sociological
Review} 63:309-330.

\_\_\_\_\_\_. 1998b. "Birds of a Feather Sing Together." \emph{Social
Forces} 77:453-485.

\_\_\_\_\_\_. 2003. "Culture and Competition: Homophily and Distancing
Explanations for Cultural Niches." \emph{American Sociological Review}
68:319-345.

McPherson, J. Miller, Pamela A. Popielarz, and Sonja Drobnic. 1992.
"Social Networks and Organizational Dynamics." \emph{American
Sociological Review} 57:153-170.

McPherson, Miller. 1983. "An Ecology of Affiliation." \emph{American
Sociological Review} 48:519-532.

\_\_\_\_\_\_. 2004. "A Blau Space Primer: Prolegomenon to an Ecology of
Affiliation." \emph{Industrial and Corporate Change} 13: 263-280.

McPherson, Miller and Lynn Smith-Lovin. 1987. "Homophily in Voluntary
Organizations: Status Distance and the Composition of Face-to-Face
Groups." \emph{American Sociological Review} 52:370-379.

McPherson, Miller, Lynn Smith-Lovin, and James M. Cook. 2001. "Birds of
a Feather: Homophily in Social Networks." \emph{Annual Review of
Sociology} 27:415-444.

Milardo, Robert M. 1987. "Changes In Social Networks Of Women And Men
Following Divorce: A Review." \emph{Journal of Family Issues} 8: 78-96.

Mische, Ann. 2003. "Cross-talk in Movements: Reconceiving the
Culture-Network Link." Pp. 258-280 in \emph{Social Movements and
Networks: Relational Approaches to Collective Action}, edited by M.
Diani and D. McAdam. London: Oxford University Press.

Morgan, David, Paula Carder, and Margaret Neal. 1997. "Are Some
Relationships more Useful than Others? The Value of Similar Others in
the Networks of Recent Widows." \emph{Journal of Social and Personal
Relationships} 14:745-759.

Morgan, David L., Margaret B. Neal, and Paula Carde. 1997. "The
stability of core and peripheral networks over time." \emph{Social
Networks} 19:9-25.

Munch, Allison, J. Miller McPherson, and Lynn Smith-Lovin. 1997.
"Gender, Children, and Social Contact: The Effects of Childrearing for
Men and Women." \emph{American Sociological Review} 62:509-520.

Ostrower, Francie. 1998. "The Arts as Cultural Capital among Elites:
Bourdieu\textquotesingle s Theory Reconsidered." \emph{Poetics}
26:43-53.

Pahl, Ray and David J. Pevalin. 2005. "Between family and friends: a
longitudinal study of friendship choice." \emph{British Journal of
Sociology} 56:433--450.

Peterson, Richard A. 1992. "Understanding Audience Segmentation: From
Elite and Popular to Omnivore and Univore." \emph{Poetics} 21:243-258.

\_\_\_\_\_\_. 1997. "The Rise and Fall of Highbrow Snobbery as a Status
Marker." \emph{Poetics} 25:75-92.

\_\_\_\_\_\_. 2005. "Problems in Comparative Research: The Example of
Omnivorousness." \emph{Poetics} 33:257-282.

Popielarz, Pamela A. and J. Miller McPherson. 1995. "On the Edge or In
Between: Niche Position, Niche Overlap, and the Duration of Voluntary
Association Memberships." \emph{American Journal of Sociology}
101:698-720.

Relish, Michael. 1997. "It\textquotesingle s Not All Education: Network
Measures as Sources of Cultural Competency." \emph{Poetics} 25:121-139.

Simmel, Georg. 1949. "The Sociology of Sociability." \emph{American
Journal of Sociology} 55:254-261.

Simon, Herbert A. 1952. "Comments on the Theory of Organizations."
\emph{American Political Science Review} 46:1130-1139.

Smith, Tom W. 1995. "Generational Differences in Musical Preferences."
\emph{Popular Music and Society} 18:43-59.

Stinchcombe, Arthur L. 1965. "Social Structure and Organizations." Pp.
142-193 in \emph{Handbook of Organizations}, edited by J. G. March.
Chicago: Rand McNally.

Straits, Bruce C. 1996. "Ego-Net Diversity: Same- and Cross-Sex Coworker
Ties." \emph{Social Networks} 18:29-45.

Suitor, J. Jill, Barry Wellman, and David L. Morgan. 1997.
"It\textquotesingle s about time: How, why, and when networks change."
\emph{Social Networks} 19:1-7.

Suitor, Jill and Shirley Keeton. 1997. "Once a friend, always a friend?
Effects of homophily on women\textquotesingle s support networks across
a decade." \emph{Social Networks} 19:51-62.

van Eijck, Koen and Kees van Rees. 2000. "Media Orientation and Media
Use: Television Viewing Behavior of Specific Reader Types from 1975 to
1995." \emph{Communication Research} 27:574-616.

van Rees, Kees, Jeroen K. Vermunt, and Marc Verboord. 1999. "Cultural
Classifications Under Discussion: Latent Class Analysis of Highbrow and
Lowbrow Reading." \emph{Poetics} 26:349-365.

Wali, Alaka, Rebecca Severson, and Mario Longoni. 2002. \emph{Informal
Arts: Finding Cohesion, Capacity and Other Cultural Benefits in
Unexpected Places}. Chicago: Chicago Center for Arts Policy.

Wellman, Barry. 1979. "The Community Question: The Intimate Networks of
East Yorkers." \emph{American Journal of Sociology} 84:1201-1231.

\_\_\_\_\_\_. 1983. "Network Analysis: Some Basic Principles."
\emph{Sociological Theory} 1:155-200.

\_\_\_\_\_\_. 1988. "Structural Analysis: from Method and Metaphor to
Theory and Substance." Pp. 19-61 in \emph{Social Structures a Network
Approach}, edited by B. Wellman and S. D. Berkowitz. Cambridge:
Cambridge University Press.

Wellman, Barry and S. D. Berkowitz. 1988. "Introduction: Studying Social
Structures." Pp. 1-18 in \emph{Social Structures a Network Approach},
edited by B. Wellman and S. D. Berkowitz. Cambridge: Cambridge
University Press.

Wellman, Barry, Peter J. Carrington, and Alan Hall. 1988. "Networks as
personal communities." Pp. 130-184 in \emph{Social Structures: a Network
Approach}, edited by B. Wellman and S. D. Berkowitz. Cambridge:
Cambridge University Press.

Wellman, Barry, Renita Yuk-Lin Wong, David Tindall, and Nancy Nazer.
1997. "A Decade of Network Change: Turnover, Persistence and Stability
in Personal Communities." \emph{Social Networks} 19:27--50.

Wellman, Barry and Scot Wortley. 1990. "Different Strokes from Different
Folks: Community Ties and Social Support." \emph{American Journal of
Sociology} 96:558-588.

\includegraphics[width=6.01319in,height=3.10417in]{media/image1.emf}

Figure

Table 1.

\includegraphics[width=6.81597in,height=1.83194in]{media/image2.wmf}

Table 2.

Table 3. Lagged, Within-Between (Mundlak) Poisson Models
Effects Models of the effect of culture consumption diversity on
interaction with friends and number of organizational memberships,
Project Canada Survey 1985-1990.

\begin{longtable}[]{@{}
  >{\raggedright\arraybackslash}p{(\columnwidth - 14\tabcolsep) * \real{0.2963}}
  >{\raggedright\arraybackslash}p{(\columnwidth - 14\tabcolsep) * \real{0.1174}}
  >{\raggedright\arraybackslash}p{(\columnwidth - 14\tabcolsep) * \real{0.1174}}
  >{\raggedright\arraybackslash}p{(\columnwidth - 14\tabcolsep) * \real{0.1132}}
  >{\raggedright\arraybackslash}p{(\columnwidth - 14\tabcolsep) * \real{0.0234}}
  >{\raggedright\arraybackslash}p{(\columnwidth - 14\tabcolsep) * \real{0.1048}}
  >{\raggedright\arraybackslash}p{(\columnwidth - 14\tabcolsep) * \real{0.1084}}
  >{\raggedright\arraybackslash}p{(\columnwidth - 14\tabcolsep) * \real{0.1190}}@{}}
\toprule\noalign{}
\begin{minipage}[b]{\linewidth}\raggedright
\end{minipage} &
\multicolumn{3}{>{\raggedright\arraybackslash}p{(\columnwidth - 14\tabcolsep) * \real{0.3481} + 4\tabcolsep}}{%
\begin{minipage}[b]{\linewidth}\raggedright
Interaction with Friends
\end{minipage}} & \begin{minipage}[b]{\linewidth}\raggedright
\end{minipage} &
\multicolumn{3}{>{\raggedright\arraybackslash}p{(\columnwidth - 14\tabcolsep) * \real{0.3322} + 4\tabcolsep}@{}}{%
\begin{minipage}[b]{\linewidth}\raggedright
Number of Memberships
\end{minipage}} \\
\midrule\noalign{}
\endhead
\bottomrule\noalign{}
\endlastfoot
& Lagged & Random

Effects & IV

Random

Effects & & Lagged & Random

Effects & IV

Random

Effects \\
N. of Cultural Tastes & 0.0289* & 0.119* & 0.142* & & 0.0874* & 0.0504*
& 0.155* \\
& (3.46) & (17.15) & (5.20) & & (3.88) & (4.59) & (2.73) \\
Gender (Women=1) & -0.0623+ & -0.123* & -0.128* & & 0.284* & 0.257* &
0.116 \\
& (-1.80) & (-3.16) & (-2.50) & & (3.05) & (3.27) & (1.04) \\
Education & 0.0178 & 0.00461 & -0.00680 & & 0.0768* & 0.0907* &
0.0671* \\
& (1.60) & (0.38) & (-0.41) & & (2.58) & (3.93) & (2.06) \\
Marital Status (Married=1) & -0.0585 & -0.0610 & -0.134* & & 0.0572 &
0.00358 & 0.152 \\
& (-1.36) & (-1.31) & (-2.07) & & (0.49) & (0.04) & (1.14) \\
Children & 0.00392 & 0.00684 & 0.0155 & & 0.0173 & 0.0148 & 0.00837 \\
& (0.36) & (0.55) & (0.92) & & (0.59) & (0.59) & (0.23) \\
Age & 0.00486* & 0.000912 & 0.00144 & & 0.00194 & 0.00205 & 0.00772+ \\
& (4.35) & (0.76) & (0.78) & & (0.64) & (0.87) & (1.66) \\
Urban & -0.00485 & -0.0267 & 0.00220 & & -0.115 & -0.0743 & -0.0370 \\
& (-0.17) & (-0.82) & (0.05) & & (-1.49) & (-1.13) & (-0.39) \\
Constant & -0.310* & -0.222* & -0.310 & & -0.264 & 0.0694 & -0.777 \\
& (-3.41) & (-2.37) & (-1.56) & & (-1.08) & (0.38) & (-1.61) \\
& & & & & & & \\
R\textsuperscript{2} & 0.0766 & & & & 0.0734 & & \\
\emph{R\textsuperscript{2}} (within) & & 0.312 & 0.308 & & & 0.00704 &
0.0175 \\
\emph{R\textsuperscript{2}} (between) & & 0.123 & 0.133 & & & 0.0943 &
0.0703 \\
χ\textsuperscript{2} & & 319.7 & 51.55 & & & 54.09 & 21.85 \\
N & 560 & 1116\textsuperscript{a} & 657\textsuperscript{a, b} & & 560 &
1116\textsuperscript{a} & 657\textsuperscript{a, b} \\
\end{longtable}

+p\textless0.10; *p\textless0.05; **p\textless0.01

\textsuperscript{a}Unit of analysis is the individual/time-period. Each
individual contributes two observations for the analysis (1985, 1990).

\textsuperscript{b}Only those individuals with data for three panels
(1980, 1985, 1990) can be included in the analysis.

\section{\texorpdfstring{Appendix }{Appendix }}\label{appendix}

Table 4. Means, standard deviations and range of the variables used in
the analysis. Project Canada Survey 1985-1995

\begin{longtable}[]{@{}
  >{\raggedright\arraybackslash}p{(\columnwidth - 10\tabcolsep) * \real{0.5016}}
  >{\raggedright\arraybackslash}p{(\columnwidth - 10\tabcolsep) * \real{0.0740}}
  >{\raggedright\arraybackslash}p{(\columnwidth - 10\tabcolsep) * \real{0.1255}}
  >{\raggedright\arraybackslash}p{(\columnwidth - 10\tabcolsep) * \real{0.1356}}
  >{\raggedright\arraybackslash}p{(\columnwidth - 10\tabcolsep) * \real{0.0792}}
  >{\raggedright\arraybackslash}p{(\columnwidth - 10\tabcolsep) * \real{0.0842}}@{}}
\toprule\noalign{}
\begin{minipage}[b]{\linewidth}\raggedright
Variable
\end{minipage} & \begin{minipage}[b]{\linewidth}\raggedright
N
\end{minipage} & \begin{minipage}[b]{\linewidth}\raggedright
Mean
\end{minipage} & \begin{minipage}[b]{\linewidth}\raggedright
Std Dev.
\end{minipage} & \begin{minipage}[b]{\linewidth}\raggedright
Min
\end{minipage} & \begin{minipage}[b]{\linewidth}\raggedright
Max
\end{minipage} \\
\midrule\noalign{}
\endhead
\bottomrule\noalign{}
\endlastfoot
Δ Social Interaction & 561 & 0.81 & -\/- & 0 & 1 \\
Δ Organizational Memberships & 561 & 0.79 & -\/- & 0 & 1 \\
Δ Educational Status & 561 & 0.33 & -\/- & 0 & 1 \\
Δ Geographic Location & 561 & 0.53 & -\/- & 0 & 1 \\
Δ Work Status & 561 & 0.24 & -\/- & 0 & 1 \\
Δ Marital Status & 561 & 0.08 & -\/- & 0 & 1 \\
Δ N. of Children & 561 & 0.08 & -\/- & 0 & 1 \\
Marital Status in 1985 (Married=1) & 561 & 0.83 & -\/- & 0 & 1 \\
Marital Status in 1990 (Married=1) & 561 & 0.81 & -\/- & 0 & 1 \\
Work Status in 1985 (Employed=1) & 561 & 0.65 & -\/- & 0 & 1 \\
Work Status in 1990 (Employed=1) & 561 & 0.58 & -\/- & 0 & 1 \\
Gender (Women=1) & 561 & 0.67 & -\/- & 0 & 1 \\
City Size in 1985 (Large City=1) & 561 & 0.53 & -\/- & 0 & 1 \\
City Size in 1990 (Large City=1) & 561 & 0.52 & -\/- & 0 & 1 \\
Education in 1985 & 561 & 2.71 & 1.30 & 0 & 5 \\
Education in 1990 & 561 & 2.85 & 1.29 & 0 & 5 \\
Age in 1985 & 561 & 51.53 & 13.30 & 19 & 87 \\
Age in 1990 & 561 & 55.96 & 14.12 & 24 & 88 \\
N. of Organizational Memberships in 1985 & 561 & 1.97 & 1.40 & 0 & 6 \\
N. of Organizational Memberships in 1990 & 561 & 1.89 & 1.36 & 0 & 9 \\
Freq. of Social Interaction in 1985 & 561 & 1.54 & 0.58 & 1 & 3 \\
Freq. of Social Interaction in 1990 & 561 & 2.78 & 1.02 & 1 & 5 \\
Culture Consumption Scale in 1980 & 339 & 6.76 & 1.19 & 3 & 8 \\
Culture Consumption Scale in 1985 & 560 & 6.43 & 1.72 & 1 & 10 \\
Culture Consumption Scale in in 1990 & 556 & 3.74 & 1.40 & 0 & 8 \\
\end{longtable}

It is well known that measurement error attenuates the relationships
between variables. This means that the relationship between two observed
variables which are subject to measurement error will generally be
weaker than their true relationship. For the analysis of change, this
phenomenon implies that when the observed states are subject to
measurement error, the strength of the relationships among the true
states occupied at two subsequent points in time will be underestimated,
or in other words, the amount of change will be overestimated. When the
data are subject to measurement error, the observed transitions are, in
fact, a mixture of true change and spurious change resulting from
measurement error (Hagenaars, 1992; Van de Pol \& De Leeuw, 1986). The
latent Markov model makes it possible to separate true change and
spurious change caused by measurement error (Vermunt et al, 1999: 184).

\#

Mixed Markov Latent Class Models

\# Frank van de Pol and Rolf Langeheine

\# Sociological Methodology, Vol. 20, (1990), pp. 213-247

\end{document}
